\cotefile{14}
\scanpage{101}\mspage{\handwritten{95}}
% archive sheet number: 100
\diagram{\begin{tikzcd}[ampersand replacement=\&]
 E'_{j,n}\handwritten{\arrow[r,"{\handwritten{v_i}}"]} \& E^{\deleted{\unsure{'}}}_{i,n}\\
 E'_{k,n}\handwritten{\arrow[u]\arrow[r,no head,"{\handwritten{v_j}}"]} \& E_{j,n}\handwritten{\arrow[u]}
\end{tikzcd}}
où les flêches verticales sont les homomorphismes de transition.

N.B. On pose $E_{i,n}=E_i/I_i^{n+1}$, $E'_{i,n}=\operatorname{Symm}_{\unsure{C|O}_i}(I_i/I_i^2)/$puissance $(N+1)$ème de l'\unsure{idéah|idéal} d'augmentation. La démonstration du lemme B est immédiate, en procédant comme dans le cas particulier 2.2. déjà traité.

\uline{Lemme} C \unsure{Supposons} $C$ et $B$ \unsure{form} simple sur $A$, (donc en vertu de lemme A les conditions de Lemme B sont vérifiées). Alors:

\handwritten{\unsure{a)}} Pour tout \deleted{\unsure{$i$, existe un $j>i$ et un $v_i$ comme dans lemme B, a)}, avec \uline{en plus} un} système $i,j$, $v_i$ comme dans Lemme B,\inserted{ a)} \deleted{il existe un}\inserted{\handwritten{il existe}} $k\handwritten{\geq}j$ \inserted{\handwritten{et}}\deleted{\unsure{tel que} rendant valable la conclusion de lemme B, b), et tel qu'il existe} un homomorphisme de $A$-algèbres $w_j:E_{k,n}\handwritten{\to}E'_{j,n}$ ayant les propriétés suivantes: 1º) $\handwritten{w_k}$ induit un épimorphisme sur les gradués associés \quad 2º) $v_iw_j$ est l'homomorphisme de transition $E_k\handwritten{\to}E_{\unsure{i}}$

\deleted{b) Soient $i,j,k$, $v_i,v_j$, $w_j$ comme dans a). Alors il existe un $\handwritten{\unsure{l}}\quad k$, un homomorphisme $v_k:E'_{\handwritten{\unsure{l}},n}\quad E_{k,n}$ et un homomorphisme $w_k:E_{\handwritten{\unsure{l}},n}\ \unsure{E}\ E'_{k,n}$ ayant les propriétés suivantes: 1º) $v_k$ a les propriétés mentionnées déjà pour $v_i,v_j$ \quad 2º) $w_k$ a les deux propriétés mentionnées déjà pour $w_j$ \quad \unsure{3}º}
